El equipo utilizó una herramienta de inteligencia artificial generativa como apoyo durante la elaboración de esta entrega. La Tabla~\ref{tab:ia} detalla para qué se usó, en qué partes del trabajo y cómo se verificó el resultado.

\begin{table}[H]
  \centering\small
  \caption{Uso de inteligencia artificial en la entrega TA-PFC-E1.}\label{tab:ia}
  \renewcommand{\arraystretch}{1.3}
  \begin{tabularx}{\textwidth}{@{}>{\raggedright\arraybackslash}p{2.4cm}LL>{\raggedright\arraybackslash}p{4.2cm}@{}}
    \toprule
    Herramienta & Propósito & Parte del trabajo & Cómo se verificó \\
    \midrule
    Claude Opus 5.5 (Anthropic) & Proponer la estructura de carpetas del repositorio y la plantilla del documento principal (\texttt{main.tex}) & Estructura del repositorio y del documento & Se comparó con la estructura exigida en la guía de la tarea y se compiló siguiendo el \texttt{README.md}. \\
    Claude Opus 5.5 (Anthropic) & Revisar la redacción de los instrumentos elaborados por el equipo y sugerir correcciones (preguntas dobles, preguntas inductivas, opciones incompletas) & Anexos A, B y C & El equipo revisó cada cambio sugerido y lo aceptó o descartó según la realidad del gimnasio observada el 18/11/2026. \\
    Claude Opus 5.5 (Anthropic) & Dar formato LaTeX a la guía de entrevista, al cuestionario y al consentimiento informado redactados por el equipo & Anexos A, B y D & Se compararon los textos compilados con los documentos originales del equipo. \\
    Claude Opus 5.5 (Anthropic) & Proponer un borrador de las tablas de roles, cronograma y plan de comunicación & Sección 7 & El equipo revisó los roles según el reparto acordado y ajustó las fechas a las sesiones de la Tabla~\ref{tab:sesiones}. \\
    Claude Opus 5.5 (Anthropic) & Revisar la coherencia entre secciones (códigos de partes interesadas, horarios, referencias cruzadas y citas) & Documento completo & Se corrigieron las diferencias detectadas y se volvió a compilar el documento sin errores. \\
    \bottomrule
  \end{tabularx}
\end{table}

Los datos de campo (visita a la Unidad de Bienestar Universitario, autorización de la organización y observación del 18/11/2026), la identificación de las partes interesadas, los objetivos de negocio y las decisiones de alcance provienen del trabajo del equipo. Ningún dato de campo, evidencia ni referencia fue generado por la herramienta. Cada integrante revisó su sección y puede explicar su contenido.
